\chapter{I/O Runtime}
\label{chap:io-runtime}

\section{Readiness}

The readiness path is small. A future registers interest in a file descriptor,
the scheduler records the waker, the Unix reactor sleeps using Linux
\texttt{epoll}, macOS/iOS \texttt{kqueue}, or fallback \texttt{poll} plus a
wake pipe, and readiness wakes the interested task.

This path is readiness-based, not completion-based. It is used by the current
TCP helpers and remains separate from the experimental \texttt{io\_uring}
completion path.

Readiness is the first I/O layer because it matches Unix non-blocking sockets
and is portable. Ownership stays simple: user code owns the stream or listener;
the runtime only tracks which task to wake when the descriptor may progress.

Readiness means the OS reports an operation is likely to make progress, not
that bytes are available. A readable stream may return fewer bytes than
requested. A writable stream may accept only part of a buffer. The future must
attempt the operation, handle \texttt{WouldBlock}, and re-register interest.

\begin{lstlisting}[language={},caption={Readiness future shape}]
poll future:
    try non-blocking operation
    if operation succeeds:
        return Ready(result)
    if operation would block:
        register fd interest and waker
        return Pending

reactor wait:
    OS reports fd readiness
    scheduler wakes interested task
\end{lstlisting}

Readiness futures do not own file descriptors. The caller owns the stream or
listener; the future borrows it for the async operation and registers interest
while pending.

A connection handler owns a stream and calls helpers against it in sequence.
The reactor is the waiting mechanism, not the connection owner.

\begin{notabene}[Readiness lesson]
Readiness is a notification that retrying may be useful. Completion is a
notification that the operation has already finished. The code keeps those
concepts separate because the ownership and cancellation rules are different.
\end{notabene}

\section{TCP helpers}

The TCP helpers are built on top of readiness futures rather than hiding a
separate runtime. They demonstrate the normal async server shape:

\begin{enumerate}
    \item put the listener or stream in non-blocking mode;
    \item accept or connect using readiness futures;
    \item spawn one handler task per accepted connection;
    \item use \texttt{TaskScope} or join handles to wait for handlers;
    \item propagate handler errors as ordinary \texttt{io::Result} values.
\end{enumerate}

Grouped TCP helper variants classify accepted handler tasks into scheduling
groups. The accept loop remains in the caller's task; only per-connection
handler work enters the requested group. This separates foreground request
handling from background maintenance or low-priority connections.

Accepting connections and serving connections have different lifecycle and
scheduling needs. The accept loop is one long-lived task. Handlers are many
shorter-lived tasks. Keeping those roles separate makes shutdown,
observability, and scheduling group assignment straightforward.

\section{Sharded TCP socket options}

\texttt{ShardedTcpServer} keeps Linux socket placement explicit. On Linux,
\texttt{SO\_REUSEPORT} lets each shard bind its own listener so the kernel can
distribute incoming connections across shard-local accept loops. The server can
also opt into \texttt{SO\_INCOMING\_CPU} with
\texttt{ShardedTcpIncomingCpu}. This option is disabled by default; callers
must request either sequential placement over \texttt{available\_cpu\_ids()} or
an explicit CPU list.

Accepted sharded TCP connections expose a Linux \texttt{kTLS ULP} helper. That helper
attaches the kernel \texttt{tls} upper-layer protocol to the socket. It is not
a complete TLS implementation: the handler or a TLS stack must still complete
the handshake and install the required kernel TX/RX crypto state. The server
does not own certificates, handshake transcripts, traffic secrets, or TLS
record policy, so it deliberately leaves the full kTLS handoff outside the
generic accept loop.

\section{\texttt{io\_uring}}

The Linux \texttt{io\_uring} backend is supported but deliberately narrow. It
integrates owned-buffer \texttt{read\_at},
\texttt{read\_exact\_at}, and \texttt{write\_at} style file I/O.

Unlike readiness, completion I/O can outlive the Rust future that submitted
it. The runtime must therefore track operation ids, owned buffers, late
completions, and abandonment.

Linux hosts and containers may restrict or disable \texttt{io\_uring}. Sitas
does not manage host policy. Setup rejection is reported as unavailable;
submission rejection is returned as an I/O error.

Each Linux executor run loop installs a thread-local dispatcher when the host
allows ring setup. Executor-backed \texttt{io\_uring} futures store operation
ids and owned buffers. When polled on a shard, they queue an operation against
the thread-local dispatcher and register their task waker.

When the executor becomes idle, it submits pending ring entries and polls the
ring descriptor through the same reactor wait used for fd readiness and timer
deadlines. Ring descriptor readiness is only a notification that completion
queue entries may be harvested. The completed operation still flows through the
dispatcher, which matches operation ids, returns owned buffers, and wakes
waiting futures.

The dispatcher tracks:

\begin{itemize}
    \item queued submissions;
    \item locally buffered completions;
    \item registered wakers;
    \item abandoned operations;
    \item deferred owned buffers;
    \item cumulative dispatched, buffered, woken, and discarded operation
          counters.
\end{itemize}

\section{Completion lifecycle}

The \texttt{io\_uring} path is completion-based. A future submits an operation
and later receives a completion. That changes the ownership problem. With
readiness, the future borrows a stream and retries ordinary system calls. With
\texttt{io\_uring}, the kernel may hold a pointer to an owned buffer until the
completion arrives.

The runtime tracks operation state explicitly:

\begin{lstlisting}[language={},caption={Simplified completion lifecycle}]
future first poll
    allocate operation id
    move owned buffer into dispatcher tracking
    submit operation to ring
    register task waker
    return Pending

completion arrives
    record result and completed buffer
    wake registered task

future next poll
    take completion
    return Ready(result, buffer)
\end{lstlisting}

If the future is dropped before completion, the operation may still be in the
kernel. The dispatcher marks it abandoned and defers buffer cleanup until a
completion or safe discard point. Cancellation therefore includes
kernel-owned in-flight work, not only waker removal.

The dispatcher holds buffers, matches completions to futures, wakes tasks, and
accounts for futures dropped before the kernel replied.

\section{Current limitations}

Timers, readiness waits, and ring completions share one Linux executor idle
wait shape, but they are not exposed as one portable event source. Completion
dispatch uses a fixed per-tick budget. Executor snapshots expose the lifecycle
state, dispatch counts, budget exhaustion, and the final shutdown-drain result.

Setting \texttt{SITAS\_REQUIRE\_IO\_URING=1} converts an unavailable backend
into an executor setup error. During teardown, the dispatcher makes a bounded
attempt to drain abandoned work when no live task wakers remain. If
\texttt{run\_until} returns while spawned ring work is pending, the dispatcher
stays installed on that executor thread so a later run can continue it.

The runtime does not yet provide a portable completion source, caller-tunable
completion policy, or a production shutdown policy for mixed readiness and
completion workloads.

\section{Running statistics}

\href{https://en.wikipedia.org/wiki/Algorithms_for_calculating_variance\#Welford's_online_algorithm}{Welford's
algorithm} computes mean, variance, and standard deviation in a single pass
without storing individual samples. The \texttt{running\_stats} module provides
\texttt{RunningStatistics} as a small online accumulator useful for runtime
diagnostics and benchmark comparisons.

\begin{lstlisting}[language={},caption={RunningStatistics use}]
let mut stats = RunningStatistics::new();
stats.add_sample_duration(elapsed);       // record as seconds (f64)
stats.add_sample_interval(start, end);    // convenience for Instant pairs
stats.merge(&other);                      // Chan et al. parallel merge

// Accessors return Option<f64> (None until enough samples):
stats.count()         // u64
stats.min() / max()   // Option<f64>
stats.mean()          // Option<f64>
stats.variance()      // Option<f64>,  sample variance (n-1), n >= 2
stats.std_dev()       // Option<f64>
stats.cv()            // Option<f64>,  coefficient of variation in %
stats.total()         // u64 integral total
\end{lstlisting}

Duration-valued accessors are available for the mean, minimum, maximum, and
standard deviation; they return \texttt{Option<Duration>} instead of seconds.

The \texttt{merge} method combines per-shard accumulators with Chan et al.'s
parallel algorithm without revisiting the original samples.

The display format prints count, mean, standard deviation, CV, min, max, and
total in a single human-readable line, which is how the index build examples
report per-shard timing statistics.

CV is unit-free, so it can compare dispersion across differently scaled
measurements.
